% GENERATED FILE. Edit canonical lessons, then run npm run generate:books.
% canonical-source-hash: fnv1a64:03349651b00732d7
% canonical-lessons: BN-C117-shim, BN-C117-kumro, BN-C117-shak, BN-C117-mulo, BN-C117-chal

\chapter{Beans, Pumpkin, Leafy greens, Radish, Uncooked rice}
\label{ch:bn-beans-pumpkin-leafy-greens-radish-uncooked-rice}
\hlchaptermodality{117}

\begin{chapteropening}
\textbf{By the end of this chapter:} \emph{I can say beans, a pumpkin, leafy greens, a radish and uncooked rice.}

Complete the last lesson of chapter 117: i can say beans, a pumpkin, leafy greens, a radish and uncooked rice.
\end{chapteropening}

\section[\texorpdfstring{\bn{শিম} (shim)}{শিম (shim)}]{\bn{শিম} (shim) — beans}

\label{lesson:BN-C117-shim}



\begin{quote}
Before the new one: say the Bengali for ginger, then the Bengali for a tomato.
\end{quote}

\subsection*{You'll want to know: \bn{শিম}}
\textbf{\bn{শিম}} — \emph{shim} — \textquotedblleft{}beans\textquotedblright{}.

\begin{grammarlens}[title={What you've built}]
One more word for this chapter.
\end{grammarlens}

\subsection*{Your turn}
\begin{itemize}
\raggedright
  \item \emph{Say it:} \emph{shim}
  \item \emph{Say it:} \emph{shim}, once more
  \item \emph{Say it:} say \emph{shim}
  \item \emph{Recall:} say the Bengali for ginger, then the Bengali for a tomato, then say \emph{shim} again
\end{itemize}

\begin{tcolorbox}[breakable,colback=teal!4,colframe=teal!35!black,title={Before you move on}]
What does \bn{শিম} mean? (\textquotedblleft{}Beans\textquotedblright{}.) Say it once more.
\end{tcolorbox}
\section[\texorpdfstring{\bn{কুমড়ো} (kumṛo)}{কুমড়ো (kumṛo)}]{\bn{কুমড়ো} (kumṛo) — a pumpkin}

\label{lesson:BN-C117-kumro}



\begin{quote}
Before the new one: say the Bengali for a tomato, then the Bengali for beans.
\end{quote}

\subsection*{You'll want to know: \bn{কুমড়ো}}
\textbf{\bn{কুমড়ো}} — \emph{kumṛo} — \textquotedblleft{}a pumpkin\textquotedblright{}.

\begin{grammarlens}[title={What you've built}]
One more word for this chapter.
\end{grammarlens}

\subsection*{Your turn}
\begin{itemize}
\raggedright
  \item \emph{Say it:} \emph{kumṛo}
  \item \emph{Say it:} \emph{kumṛo}, once more
  \item \emph{Say it:} say \emph{kumṛo}
  \item \emph{Recall:} say the Bengali for a tomato, then the Bengali for beans, then say \emph{kumṛo} again
\end{itemize}

\begin{tcolorbox}[breakable,colback=teal!4,colframe=teal!35!black,title={Before you move on}]
What does \bn{কুমড়ো} mean? (\textquotedblleft{}A pumpkin\textquotedblright{}.) Say it once more.
\end{tcolorbox}
\section[\texorpdfstring{\bn{শাক} (shāk)}{শাক (shāk)}]{\bn{শাক} (shāk) — leafy greens}

\label{lesson:BN-C117-shak}



\begin{quote}
Before the new one: say the Bengali for beans, then the Bengali for a pumpkin.
\end{quote}

\subsection*{You'll want to know: \bn{শাক}}
\textbf{\bn{শাক}} — \emph{shāk} — \textquotedblleft{}leafy greens\textquotedblright{}.

\begin{grammarlens}[title={What you've built}]
One more word for this chapter.
\end{grammarlens}

\subsection*{Your turn}
\begin{itemize}
\raggedright
  \item \emph{Say it:} \emph{shāk}
  \item \emph{Say it:} \emph{shāk}, once more
  \item \emph{Say it:} say \emph{shāk}
  \item \emph{Recall:} say the Bengali for beans, then the Bengali for a pumpkin, then say \emph{shāk} again
\end{itemize}

\begin{tcolorbox}[breakable,colback=teal!4,colframe=teal!35!black,title={Before you move on}]
What does \bn{শাক} mean? (\textquotedblleft{}Leafy greens\textquotedblright{}.) Say it once more.
\end{tcolorbox}
\section[\texorpdfstring{\bn{মুলো} (mulo)}{মুলো (mulo)}]{\bn{মুলো} (mulo) — a radish}

\label{lesson:BN-C117-mulo}



\begin{quote}
Before the new one: say the Bengali for a pumpkin, then the Bengali for leafy greens.
\end{quote}

\subsection*{You'll want to know: \bn{মুলো}}
\textbf{\bn{মুলো}} — \emph{mulo} — \textquotedblleft{}a radish\textquotedblright{}.

\begin{grammarlens}[title={What you've built}]
One more word for this chapter.
\end{grammarlens}

\subsection*{Your turn}
\begin{itemize}
\raggedright
  \item \emph{Say it:} \emph{mulo}
  \item \emph{Say it:} \emph{mulo}, once more
  \item \emph{Say it:} say \emph{mulo}
  \item \emph{Recall:} say the Bengali for a pumpkin, then the Bengali for leafy greens, then say \emph{mulo} again
\end{itemize}

\begin{tcolorbox}[breakable,colback=teal!4,colframe=teal!35!black,title={Before you move on}]
What does \bn{মুলো} mean? (\textquotedblleft{}A radish\textquotedblright{}.) Say it once more.
\end{tcolorbox}
\section[\texorpdfstring{\bn{চাল} (chāl)}{চাল (chāl)}]{\bn{চাল} (chāl) — uncooked rice}

\label{lesson:BN-C117-chal}



\begin{quote}
Before the new one: say the Bengali for leafy greens, then the Bengali for a radish.
\end{quote}

\subsection*{You'll want to know: \bn{চাল}}
\textbf{\bn{চাল}} — \emph{chāl} — \textquotedblleft{}uncooked rice\textquotedblright{}.

\begin{grammarlens}[title={What you've built}]
One more word for this chapter.
\end{grammarlens}

\subsection*{Your turn}
\begin{itemize}
\raggedright
  \item \emph{Say it:} \emph{chāl}
  \item \emph{Say it:} \emph{chāl}, once more
  \item \emph{Say it:} say \emph{chāl}
  \item \emph{Recall:} say the Bengali for leafy greens, then the Bengali for a radish, then say \emph{chāl} again
\end{itemize}

\begin{tcolorbox}[breakable,colback=teal!4,colframe=teal!35!black,title={Before you move on}]
What does \bn{চাল} mean? (\textquotedblleft{}Uncooked rice\textquotedblright{}.) Say it once more.
\end{tcolorbox}
